%======================================================================
% L03a_Entropy.pdf
%======================================================================
\providecommand{\dbar}{{\mathchar'26\mkern-12mu d}}
\providecommand{\modreB}[1]{{\color{NavyBlue}\small\bfseries #1}}
\providecommand{\cancel}[1]{\text{\sout{$\displaystyle #1$}}}

\section{Shrnutí -- druhý TD zákon, absolutní teplota}
\strana{L03a-1}

\begin{veta}[2.\ TD zákon]
(Clausius) Pokud teplo samovolně přechází ze syst.\ $A$ do systému $B$, potom není možné realizovat proces, jehož jediným výsledkem by byl přenos tepla z $B$ do $A$.
\[ \Updownarrow \]
(Kelvin) Není možné realizovat proces, jehož jediným výsledkem by bylo získání tepla z lázně a jeho transformace na práci.
\end{veta}

\underline{Důsledky:}
\begin{enumerate}[label=\arabic*)]
  \item \underline{Carnotův teorém}: Vratný stroj je ze všech strojů pracujících mezi stejnými teplotami ten nejúčinnější.
  \item[1b)] $\forall$ vratné stroje pracující mezi stejnými lázněmi mají stejnou účinnost \ $\eta_R = \eta_R(\vartheta_H,\vartheta_C)$
  \item Carnotův cyklus -- realizace takového vratného stroje
  \item lze zavést \underline{absolutní TD teplotu} $T$: \quad $\vzorec{\eta_R = 1 - \dfrac{T_C}{T_H}}$
  \[
    \Rightarrow\ (\text{Kelvinova stupnice}) \quad T = 273{,}16\,\big(1 - \eta_R(\vartheta\sim T,\ \vartheta_0\sim 273{,}16)\big)\ \text{K}
  \]
  $\Rightarrow$ $T \geq 0$ \\
  $\Rightarrow$ $\exists$ jednoznačně def.\ \kw{absolutní nula} -- teplota lázně, ze které již není možné odebrat žádné teplo
  \item \underline{Clausiova nerovnost}
  \begin{itemize}
    \item obecný cyklus mezi $T_H$ a $T_C$ $\to$ $\eta \equiv 1 + \dfrac{Q_C}{Q_H} \leq 1 - \dfrac{T_C}{T_H}$
          \quad $\Rightarrow$ \ $\dfrac{Q_H}{T_H} + \dfrac{Q_C}{T_C} \leq 0$
    \item obecný cyklus se spojitým průběhem teploty okolí $T'$:
    \[
      \vzorec{\oint\frac{\dbar Q}{T'} \leq 0}
      \qquad \text{\parbox{0.5\textwidth}{$\to$ $\leq 0$ značí, že $\frac{\dbar Q}{T'}$ je v abs.\ hodnotě větší, pokud teplo odtéká do okolí než když přitéká do systému}}
    \]
    \item vratný cyklus
    \[
      \vzorec{\oint_{(rev)}\frac{\dbar Q}{T} = 0}
    \]
  \end{itemize}
\end{enumerate}

%----------------------------------------------------------------------
\strana{L03a-2}
\begin{itemize}
  \item \underline{pozorování}: \ $\dbar Q > 0$ $\Rightarrow$ teplo teče \underline{do} systému $\Leftrightarrow$ $T_S < T$
  \[
    \Rightarrow\ \frac{\dbar Q}{T_S} \geq \frac{\dbar Q}{T}
    \qquad {\color{NavyBlue}\frac{Q_H}{T_H} \leq -\frac{Q_C}{T_C} \ \Rightarrow\ \fcolorbox{NavyBlue}{white}{$\dfrac{Q_H}{T_H} \leq \dfrac{|Q_C|}{T_C}$}}
  \]
  \hspace*{1.5em}$\dbar Q < 0$ $\Rightarrow$ teplo teče \underline{ze} systému \quad $T_S > T$
  \[ \Rightarrow\ \frac{|\dbar Q|}{T_S} \leq \frac{|\dbar Q|}{T} \]
  $\Rightarrow$ veličina $\dfrac{|\dbar Q|}{T}$ je v abs.\ hodnotě vždy větší v \uv{přijímajícím} systému \\
  $\Rightarrow$ Clausiova nerovnost je odrazem tohoto přirozeného toku tepla
\end{itemize}

\[ \Longrightarrow \quad \text{\Large\color{hlavni}\underline{ENTROPIE}} \]

\begin{itemize}
  \item omezíme-li se na vratné procesy, potom $\displaystyle\int_{A(rev)}^{B}\frac{\dbar Q}{T}$ \underline{nezávisí na trajektorii} ve stav.\ prostoru \\
  $\Rightarrow$ \underline{definuje stavovou veličinu} \quad $S = S(\xi_1,\dots,\xi_n)$: \quad {\color{NavyBlue}$\xi$: $U, X_1, \dots, X_n$}
\end{itemize}

\begin{minipage}{0.42\textwidth}
\begin{tikzpicture}[>=Stealth,scale=0.85]
  \draw[->] (-0.2,0) -- (5.4,0) node[below] {$V$};
  \draw[->] (0,-0.2) -- (0,3.6);
  \draw (0.8,0.6) .. controls (0.2,1.2) and (0.4,2.4) .. (1.2,2.8)
                 .. controls (2,3.2) and (2.8,2.8) .. (3.4,3.1)
                 .. controls (3.8,3.2) and (3.6,2.6) .. (3.3,2.2)
                 .. controls (3,1.6) and (3.4,0.8) .. (2.8,0.6)
                 .. controls (2.2,0.4) and (1.4,0.4) .. (0.8,0.6);
  \fill (0.8,0.6) circle (1.8pt) node[below left] {$O$};
  \fill (3.55,2.85) circle (1.8pt) node[above right] {$A$};
  \draw[->] (1.9,3.02) -- (2.1,3.0); \node at (1.8,3.35) {$I$};
  \draw[->] (3.32,1.5) -- (3.35,1.75); \node at (3.7,1.6) {$II$};
\end{tikzpicture}
\end{minipage}
\begin{minipage}{0.56\textwidth}
\[
  \int_{O(I)}^{A}\frac{\dbar Q}{T} - \int_{O(II)}^{A}\frac{\dbar Q}{T} = 0
  \ \Rightarrow\ \int_{O(I)}^{A}\frac{\dbar Q}{T} = \int_{O(II)}^{A}\frac{\dbar Q}{T}
\]
\begin{definice}[Entropie]
\[
  \Rightarrow\ \vzorec{S(A) \equiv \int_{O(rev)}^{A}\frac{\dbar Q}{T}}
\]
\end{definice}
\end{minipage}

\begin{itemize}
  \item $S$ def.\ až na aditivní konstantu \quad $S(O)$ \dots\ entropie ref.\ stavu (stejně jako pro $U$)

  \underline{Pozn:} $\bullet$ tato aditivní konstanta může být fyzikálně významná pro procesy zahrnující látkovou výměnu při $T\to 0$ \\
  \hspace*{2.5em}$\Rightarrow$ motivace pro 3.\ TDZ (chování entropie pro $T\to 0$) \\
  \hspace*{2.5em}(viz Luscombe str.\ 125)
  \item my se ovšem zatím stále omezujeme na \underline{uzavřené systémy} \pozor{!}
  \begin{itemize}[nosep]
    \item[--] celá konstrukce od začátku stojí na možnosti systém adiabaticky izolovat
    \item[--] Clausiova nerovnost nezahrnuje výměnu částic, pouze tepla; s tokem částic je nevyhnutelně spojen také přenos energie, který nelze spojit s $\dbar Q$ \dots
  \end{itemize}
\end{itemize}

%----------------------------------------------------------------------
\strana{L03a-3}
\begin{itemize}
  \item $\dd S = \dfrac{\dbar Q}{T}$ je úplný diferenciál \& $T > 0$ \\
  $\Rightarrow$ můžeme přeformulovat 1.\ TDZ
\end{itemize}

\begin{minipage}{0.4\textwidth}
\begin{veta}[1.\ TDZ pomocí entropie]
\[ \vzorec{\dd U = T\,\dd S + \sum_j y_j\,\dd X_j} \]
\end{veta}
\end{minipage}
\hfill
\begin{minipage}{0.56\textwidth}
{\color{Red}\Large!}\ $\to$ stavový prostor je definovaný \uv{deformačními} extenzivními proměnnými $\{X_i\}_{i=1}^{N}$, spojenými s konáním práce v kvazistac.\ procesech, a jednou \uv{\underline{nedeformační}/tepelnou};
\end{minipage}

\hspace*{1.5em}$\to$ nedeformační souřadnice se mění při přenosu tepla nebo některých způsobech konání práce v nekvazistat.\ procesech (např.\ míchání)

\medskip
$\bullet$ \underline{Pozn:} \quad $\dfrac{T_i}{T_j} = -\dfrac{Q_i}{Q_j}$ \quad $\Rightarrow$ přirozený rozměr absolutní teploty je energie: \ $[T] = [U]$ \\
\hspace*{9em}$\Rightarrow$ entropie je bezrozměrná veličina \\
\hspace*{9em}$\Leftrightarrow$ $[k_B] = \mathrm{J\,K^{-1}}$ je bezrozměrná konstanta

\begin{center}
\fcolorbox{gray}{white}{\parbox{0.85\textwidth}{\small
A system is to be taken from a particular initial state to a given final state while exchanging heat and work with reversible heat and work sources. To describe a particular process it is not sufficient merely to describe the path of the system in its thermodynamic configuration space. \emph{The critical features of the process concern the manner in which the extracted heat and work are conveyed to the reversible heat and work sources.} For that purpose auxiliary systems may be employed. The auxiliary systems are the \uv{tool} or \uv{devices} used to accomplish the task at hand, or, in a common terminology, they constitute the \emph{physical engines} by which the process is effected. \hfill \modreB{Callen str.\ 118}}}
\end{center}

\subsection*{Entropie a nevratné procesy}
\begin{itemize}
  \item $\dd S$ je úplný diferenciál $\Rightarrow$ i pro nevratný proces $A\to B$ platí
  \[
    \Delta S = S(B) - S(A) = \int_{A(rev)}^{B}\frac{\dbar Q}{T},
  \]
  kde integrál běží po libovolné vratné trajektorii spojující $A$ a $B$
\end{itemize}

{\color{Red}\small Je zřejmé, že samotná kvazistatičnost procesu nezajišťuje vratnost; důležitá je také \uv{(ne)rovnovážnost interakce s okolím}; není-li systém v rovnováze s okolím, potom musí interakce probíhat skrze vratné motory/čerpadla, která zajistí vratnou výměnu energie mezi systémem, okolím a vratnými rezervoáry tepla a práce.}

\begin{minipage}{0.36\textwidth}
\begin{tikzpicture}[>=Stealth,scale=0.85]
  \fill (0,0) circle (1.8pt) node[below] {$A$};
  \fill (3.6,2.6) circle (1.8pt) node[right] {$B$};
  \draw (0,0) .. controls (0.6,0.6) and (0,2) .. (1,2.5)
              .. controls (1.8,2.9) and (2.8,2.5) .. (3.6,2.6);
  \draw[->] (0.62,2.08) -- (0.75,2.25);
  \node[align=center] at (-0.4,2.3) {$II$\\ \small(vratný)};
  \draw[dotted,thick] (0,0) .. controls (2,0.1) and (3.3,1.2) .. (3.6,2.6);
  \draw[->] (3.05,1.25) -- (3.15,1.42);
  \node[align=center] at (3.6,0.8) {$I$\\ \small(nevratný)};
\end{tikzpicture}
\end{minipage}
\begin{minipage}{0.62\textwidth}
\[
  \Rightarrow\ \oint\frac{\dbar Q}{T'} = \int_I\frac{\dbar Q}{T'} - \int_{II}\frac{\dbar Q}{T}\leq 0
  \qquad {\color{NavyBlue}\text{\uv{$T'$}!}}
\]
\begin{veta}[Entropie v nevratném procesu]
\[
  \Rightarrow\ \vzorec{S(B) - S(A) \geq \int_I\frac{\dbar Q}{T'}} \quad (II)
  \qquad {\color{NavyBlue}\to \text{úbytek } S \text{ v okolí?}}
\]
\end{veta}
\end{minipage}

\pozor{!}\ použili jsme (CN) $\Rightarrow$ $T'$ vpravo je teplota okolí! \hfill
{\color{Red}\parbox{0.4\textwidth}{! lib.\ kvazistat.\ proces ale může být realizován jako vratný, pokud je syst.\ v jeho průběhu v rovnováze s okolím!}}

$\Rightarrow$ vidíme, že v nevratném procesu je změna entropie větší, než odpovídá jejímu prostému toku z okolí; \pozor{entropie se tedy nezachovává!}

%----------------------------------------------------------------------
\strana{L03a-4}
\hfill{\color{NavyBlue}\scriptsize\parbox{0.3\textwidth}{NB: \uv{víme/postulujeme}, že alespoň v jednom směru takový proces mezi libovol.\ dvěma body existuje\dots\ oběma zjevně nemusí}}

\begin{itemize}
  \item spec.\ nevratný proces $I$ nechť je adiabatický
\end{itemize}
\begin{veta}[Princip růstu entropie]
\[
  \to\ \vzorec{\Delta S \geq 0} \quad\Leftrightarrow\quad \text{entropie izolovaného systému nikdy neklesá} \qquad {\color{Red}\underline{šipka času}}
\]
\end{veta}
\begin{itemize}[nosep]
  \item[--] pokud je adiabatický proces $I$ navíc také vratný, potom $\Delta S = 0$ $\Rightarrow$ \uline{stavy spojené vratným adiabat.\ procesem mají stejnou entropii}
\end{itemize}

\begin{itemize}
  \item nechť $A$ a $B$ v (II) jsou infinitezimálně blízké stavy (kvazistat.\ proces)
  \[
    \Rightarrow\ \vzorec{\dd S \geq \frac{\dbar Q}{T'}} \quad \text{(CN) v dif.\ tvaru}
  \]
  \begin{itemize}[nosep]
    \item $\dfrac{\dbar Q}{T'}$ je úbytek entropie okolí \hfill
          {\color{NavyBlue}\textcircled{$*$} \ $\dfrac{\dbar Q}{T} = \dfrac{\dbar Q}{T'} + \dd S_i$, \quad $\dd S_i \geq 0$}
    \item $\dd S$ je nárůst entropie systému
  \end{itemize}
  \[
    \Rightarrow\ \boxed{\dd S = \frac{\dbar Q}{T'} + \dd S_i}\ {\color{NavyBlue}\textcircled{$*$}} \quad \& \quad \dd S_i \geq 0
  \]
  \item zatímco první člen odpovídá \uv{toku entropie} z okolí, $\dd S_i$ je \underline{přírůstek entropie} v důsledku \underline{vnitřních nevratných procesů v systému} \\
  {\color{NavyBlue}\small (typicky nějaká forma \uv{disipace} -- speciálně si to můžeme představovat jako přeměnu uspořádané energie na chaotický mikroskopický pohyb, přispívá ale např.\ i konvekce tepla -- vyrovnání teplot při přechodu do rovnováhy v rámci systému)}
\end{itemize}

\begin{priklad}[Př:]
\begin{itemize}[nosep]
  \item adiabatické míchání $\Rightarrow$ $\dd S = \dd S_i$
  \item adiabat.\ expanze do vakua $\Rightarrow$ \quad -- \,''\, --
  \item tření (disipace) -- teplo je generováno mech.\ pohybem (prací) $\Rightarrow$ nárůst entropie není kompenzován poklesem jinde \\
        {\color{NavyBlue}\small -- zde vzrůstá vnitřní energie a tím i \uv{počet} možností, jak energii rozdělit mezi částice}
  \item kvazistatický přenos tepla $\Rightarrow$ $\dd S = \dfrac{\dbar Q}{T_S} = \dfrac{\dbar Q}{T_{ext}} + \dd S_i$ \quad $(T_{ext} > T_S)$ \\
        {\color{NavyBlue}\small $\uparrow$ toto je opravdu $\dd S$ -- za vratný proces definující entropii volíme rovnovážný přenos $Q$ (tedy takový, kdy stejný objem tepla přeneseme z lázně o teplotě $T_S$ místo $T_{ext}$)}
\end{itemize}
{\color{NavyBlue}\small v těchto procesech je nárůst entropie spojen se zánikem nehomogenit $\Downarrow$ nárůst počtu kompat.\ mikrostavů (\uv{konfigurací})}
\end{priklad}

\underline{Pozn:} $\bullet$ teplo skutečně teče přirozeně z teplejšího tělesa na chladnější:

\begin{center}
\begin{tikzpicture}[>=Stealth]
  \draw (0,0) circle (0.5) node {$T_1$};
  \draw (2.4,0) circle (0.5) node {$T_2$};
  \draw[->] (0.45,0.25) .. controls (1,0.7) and (1.4,0.7) .. (1.95,0.25);
  \node at (1.2,0.85) {$\dbar Q$};
  \node[right] at (3.2,0) {$\Rightarrow\ \dfrac{\dbar Q}{T_2} - \dfrac{\dbar Q}{T_1} \geq 0 \ \Rightarrow\
    \begin{cases} T_1 > T_2 \Rightarrow \dbar Q > 0 \\ T_2 > T_1 \Rightarrow \dbar Q < 0\end{cases}$};
\end{tikzpicture}
\end{center}
\modreB{veličina $|\dbar Q|/T$ je vždy větší v přijímajícím systému}

%----------------------------------------------------------------------
\strana{L03a-5}
\subsection*{Základní důsledky existence entropie}
\begin{enumerate}[label=\arabic*)]
  \item \underline{PODMÍNKA INTEGRABILITY ENTROPIE} $\Rightarrow$ vztah mezi termickou a kalorickou stavovou rovnicí
\end{enumerate}

\begin{itemize}
  \item $\dd S$ je úplný diferenciál $\Rightarrow$ smíšené druhé derivace $S$ jsou záměnné:
  \[
    \left(\frac{\partial^2 S}{\partial\xi_i\,\partial\xi_j}\right) = \left(\frac{\partial^2 S}{\partial\xi_j\,\partial\xi_i}\right) \qquad (*)
  \]
  \item $\boxed{\dd S = \dfrac{1}{T}\dd U - \sum_i\dfrac{y_i}{T}\dd X_i}$ \ $\Leftarrow$ \ 1\&2 TDZ: $\dd U = T\,\dd S + \sum y_i\,\dd X_i$ \ {\color{NavyBlue}(+)}
  \item $U = U(T,X_1,\dots,X_n)$ \ $\Rightarrow$ \
        $\dd U = \left(\pdv{U}{T}\right)_{\{X_j\}}\dd T + \sum_i\left(\pdv{U}{X_i}\right)_{T,X_{j\neq i}}\dd X_i$ \\
  {\color{NavyBlue}\small z (+) je $U$ přirozeně funkcí $S$ \& $X_i$, ale nic nám nebrání přejít k parametrizaci $T$ \& $X_i$ -- je to stejně jednoznačné určení stavu}
\end{itemize}
\[
  \Rightarrow\ \boxed{\dd S = \frac{1}{T}\left(\pdv{U}{T}\right)_{\{X_j\}}\dd T + \sum_i\frac{1}{T}\left[\left(\pdv{U}{X_i}\right)_{T,X_{j\neq i}} - y_i\right]\dd X_i}
  \qquad {\color{NavyBlue}\text{\small\parbox{0.18\textwidth}{$\swarrow$ následně je také $S = S(T,X_1,\dots)$}}}
\]
\[
  \overset{(*)}{\Rightarrow}\ \cancel{\frac{1}{T}\left(\frac{\partial^2U}{\partial X_i\partial T}\right)_{X_{j\neq i}}}
  = \cancel{\frac{1}{T}\left(\frac{\partial^2U}{\partial T\partial X_i}\right)_{T,X_{j\neq i}}}
  - \frac{1}{T}\left(\pdv{y_i}{T}\right)_{\{X_j\}} - \frac{1}{T^2}\left[\left(\pdv{U}{X_i}\right)_{T,X_{j\neq i}} - y_i\right]
\]
\pozor{!}\ $\left(\pdv{T}{X_i}\right) = 0$ \dots\ nezávislé proměnné
\[
  \Rightarrow\ \vzorec{\left(\pdv{U}{X_i}\right)_{T,X_{j\neq i}} = -T\left(\pdv{y_i}{T}\right)_{\{X_j\}} + y_i}
\]

\begin{itemize}
  \item spec.\ \underline{1-komponentní plyn}: \quad $U = U(T,V)$, \ $y_V = -p$
  \[
    \Rightarrow\ \vzorec{\left(\pdv{U}{V}\right)_T = T\left(\pdv{p}{T}\right)_V - p}
    \qquad {\color{NavyBlue}\text{\small\parbox{0.32\textwidth}{$U_{id} = U(T)$ jsme věděli z kalorim.\ experimentů; nyní to dostáváme jako podmínku konzistence termodynamiky!}}}
  \]
  \item id.\ plyn: \quad $p = \dfrac{Nk_BT}{V}$ \ $\Rightarrow$ \ $\left(\pdv{U}{V}\right)_T = 0$ \ $\Rightarrow$ \ $U_{id} = U(T)$
\end{itemize}
$\to$ lze ukázat \quad $U_{id} = \dfrac{3}{2}Nk_BT$ \qquad ($\Leftarrow$ exp.\ nebo \underline{stat.\ fyzika}) \\
\hspace*{1.5em}(vnitřní stupně volnosti částic $\Rightarrow$ faktor $\frac32$ nahrazen jinou konstantou nebo je dokonce funkcí teploty -- např.\ $\frac32\to\frac52$ pro dvou-atom.\ id.\ plyn)

{\color{NavyBlue}\small \underline{Stat.\ fyzika}: $pV \propto NT$ platí pro $\forall$ id.\ plyny $\Rightarrow$ $U \neq U(V)$ je důsledek chybějící interakce mezi částicemi. ($\forall$ id.\ plyny $\Leftrightarrow$ nezávisle na disperzní relaci)}

